Wenn das Wasser auf Blüten und Knospen gefriert, gibt es die Erstarrungswärme ab, etwa \LfEO\ (FoTaBe S.~190). Solange noch flüssiges Wasser vorhanden ist und gefriert, bleibt die Temperatur der Pflanzenteile deshalb bei \TnO, auch wenn die Luft kälter ist. Die empfindlichen Blüten erfrieren erst bei einigen Grad unter null und werden so geschützt. Man muss dazu die ganze Nacht über weiter Wasser sprühen; die entstehende Eisschicht isoliert zusätzlich.
